\section{Construction of the auxiliary matrices and a sufficient condition for the upper bound}
\label{sec:auxiliary-matrices}
\label{sec:preliminary}

Throughout this section, let $G=(V,E(G))$ be a finite simple triangle-free
graph with $n=|V|\geq1$, and let $A=A_G$.
Fix an ordering of $V$, and index the rows and columns of $A$ by this
ordering.
For $v\in V$, its degree and the signless Laplacian are
\[
\deg_G(v)=\big|\{w\in V:\{v,w\}\in E(G)\}\big|
=\sum_{w\in V}A_{vw},
\qquad
Q_G=\operatorname{diag}_{v\in V}(\deg_G(v))+A.
\]
For $C,D\subseteq V$, write
\[
A_{C,D}=(A_{cd})_{c\in C,\,d\in D},
\]
with rows and columns in the orders inherited from $V$.
When $D=\{x\}$, write $A_{C,x}$ for $A_{C,\{x\}}$.
In particular, $A_{V,x}$ is a real column vector indexed by $V$, with
coordinate $A_{zx}$ at $z\in V$.
For a real column vector $u$ indexed by $V$, write
\[
\|u\|^2=u^{\mathsf T}u=\sum_{z\in V}u(z)^2.
\]
The notation $I$ denotes the identity matrix of the order determined
by context.
An average over $x,y$ always means independent uniform samples from $V$
with replacement:
\[
\E_{x,y}f(x,y)=\frac1{n^2}\sum_{x,y\in V}f(x,y).
\]
Thus the ordered pairs are averaged with weight $1/n^2$, including
pairs with $x=y$.

\subsection{Definition of the auxiliary matrices}
\label{subsec:auxiliary-definitions}

For every ordered pair $(x,y)\in V^2$, define the $n\times5$ real matrix
\begin{equation}\label{eq:new-U-matrix}
U^{xy}=\big(U_0^{xy}\ \ U_1^{xy}\ \ U_2^{xy}\ \ U_3^{xy}\ \ U_4^{xy}\big),
\end{equation}
whose columns are the following real column vectors indexed by $V$:
\begin{align}
U_0^{xy}
&=\left(\frac An-\frac1{50}I\right)
   (A_{V,x}-A_{V,y}),
\label{eq:new-U0}\\
U_1^{xy}
&=\frac{\deg_G(x)-\deg_G(y)}n
   (A_{V,x}+A_{V,y}),
\label{eq:new-U1}\\
U_2^{xy}
&=\operatorname{diag}_{z\in V}\left(
   \frac{\deg_G(z)}n-\frac{\deg_G(x)+\deg_G(y)}{2n}\right)
   (A_{V,x}-A_{V,y}),
\label{eq:new-U2}\\
U_3^{xy}
&=\left(\frac{A^2}{n^2}+\frac{3A}{50n}-\frac1{625}I\right)
   (A_{V,x}-A_{V,y}),
\label{eq:new-U3}\\
U_4^{xy}
&=\operatorname{diag}_{z\in V}\left(
   \frac{\deg_G(z)}n-\frac{\deg_G(x)+\deg_G(y)}{2n}\right)U_0^{xy}
   \notag\\
&\quad+\frac{(\deg_G(x)-\deg_G(y))^2}{4n^2}
   (A_{V,x}-A_{V,y}).
\label{eq:new-U4}
\end{align}
These definitions apply to all ordered pairs, including nonedges.
The factor $A_{xy}$ in the averages below makes nonedges contribute zero.

For $i,j\in\{0,1,2,3,4\}$, define
\begin{align}
\mathsf D_{ij}
&=\E_{x,y}\left[\frac{A_{xy}}n
   (U_i^{xy})^{\mathsf T}U_j^{xy}\right],
\label{eq:new-D-entry}\\
\mathsf F_{ij}
&=\E_{x,y}\left[\frac{A_{xy}}n
   (U_i^{xy})^{\mathsf T}
   \left(\frac{Q_G}n-\frac7{50}I\right)U_j^{xy}\right],
\label{eq:new-F-entry}
\end{align}
and set
\begin{equation}\label{eq:new-DF-matrices}
\mathsf D=(\mathsf D_{ij})_{i,j=0}^4,
\qquad
\mathsf F=(\mathsf F_{ij})_{i,j=0}^4.
\end{equation}
Both matrices are $5\times5$, with rows and columns indexed by
$0,1,2,3,4$ in the column order of $U^{xy}$.
Equivalently,
\begin{align}
\mathsf D
&=\E_{x,y}\left[\frac{A_{xy}}n
   (U^{xy})^{\mathsf T}U^{xy}\right],
\label{eq:new-D-matrix}\\
\mathsf F
&=\E_{x,y}\left[\frac{A_{xy}}n
   (U^{xy})^{\mathsf T}
   \left(\frac{Q_G}n-\frac7{50}I\right)U^{xy}\right].
\label{eq:new-F-matrix}
\end{align}
The matrix $Q_G$ is symmetric, so the two matrices in
\eqref{eq:new-DF-matrices} are real symmetric matrices.
The trace of a square matrix $M$ is denoted by
$\tr(M)=\sum_i M_{ii}$.

\subsection{A sufficient condition for the upper bound}
\label{subsec:sufficient-condition}

\begin{lemma}\label{lem:sufficient-condition}
If there exists a real symmetric positive definite $5\times5$ matrix
$B$ such that, for every finite simple triangle-free graph $G$ with
$n=|V(G)|\geq1$,
\[
\tr(B\mathsf F)\leq0,
\]
where $\mathsf F$ is defined from $G$ by \eqref{eq:new-F-matrix},
then \eqref{eq:main-bound} holds for every finite simple triangle-free
graph $G$ of order $n\geq1$.
\end{lemma}

\begin{proof}
Let $B$ satisfy the hypotheses, and fix a finite simple triangle-free
graph $G$ with $n\geq1$ vertices.
If $E(G)=\varnothing$, then $Q_G=0$ and
$\qmin(G)=0\leq7n/50$.
We may therefore assume $E(G)\ne\varnothing$.

We first check that $\mathsf D$ is positive semidefinite and has positive
trace.
For any real column vector $u=(u_0,\ldots,u_4)^{\mathsf T}\in\R^5$,
equations~\eqref{eq:new-U-matrix} and \eqref{eq:new-D-entry} give
\begin{align}
u^{\mathsf T}\mathsf D u
&=\sum_{i,j=0}^4u_i u_j\mathsf D_{ij}\notag\\
&=\E_{x,y}\left[\frac{A_{xy}}n
  \sum_{i,j=0}^4u_i u_j(U_i^{xy})^{\mathsf T}U_j^{xy}\right]\notag\\
&=\E_{x,y}\left[\frac{A_{xy}}n\|U^{xy}u\|^2\right]\geq0.
\label{eq:new-D-psd}
\end{align}
Here $A_{xy}/n\geq0$ for every ordered pair, so $\mathsf D$ is positive
semidefinite.

For an ordered edge $(x,y)$, triangle-freeness implies
$A_{xz}A_{yz}=0$ for every $z\in V$.
The coordinate at $z$ in \eqref{eq:new-U0} is
\begin{equation}\label{eq:new-U0-coordinate}
U_0^{xy}(z)
=\frac1n\sum_{t\in V}A_{zt}(A_{tx}-A_{ty})
 -\frac1{50}(A_{zx}-A_{zy}).
\end{equation}
In particular, using $A_{xx}=A_{yy}=0$, $A_{xy}=A_{yx}=1$, symmetry of
$A$, and $A_{zt}^2=A_{zt}$, we obtain
\begin{align}
U_0^{xy}(x)
&=\frac1n\sum_{t\in V}
  (A_{xt}^2-A_{xt}A_{yt})+\frac1{50}\notag\\
&=\frac{\deg_G(x)}n+\frac1{50}>0,
\label{eq:new-U0-at-x}\\
U_0^{xy}(y)
&=\frac1n\sum_{t\in V}
  (A_{yt}A_{xt}-A_{yt}^2)-\frac1{50}\notag\\
&=-\frac{\deg_G(y)}n-\frac1{50}<0.
\label{eq:new-U0-at-y}
\end{align}
Thus $U_0^{xy}\ne0$ for every ordered edge.
Taking the trace in \eqref{eq:new-D-matrix} yields
\begin{align}
\tr(\mathsf D)
&=\E_{x,y}\left[\frac{A_{xy}}n
  \sum_{i=0}^4\|U_i^{xy}\|^2\right]\notag\\
&\geq\E_{x,y}\left[\frac{A_{xy}}n\|U_0^{xy}\|^2\right]>0.
\label{eq:new-D-positive-trace}
\end{align}
The strict inequality holds because the edge set is nonempty and every
ordered edge has positive sampling probability $1/n^2$ and a nonzero
column $U_0^{xy}$.

By the definition of the least eigenvalue,
$Q_G-\qmin(G)I$ is positive semidefinite.
Subtracting the corresponding multiple of \eqref{eq:new-D-matrix}
from \eqref{eq:new-F-matrix}, for every $u\in\R^5$ we obtain
\begin{align}
&u^{\mathsf T}\left[
\mathsf F-\left(\frac{\qmin(G)}n-\frac7{50}\right)\mathsf D
\right]u\notag\\
&=\E_{x,y}\left[\frac{A_{xy}}n
(U^{xy}u)^{\mathsf T}
\left(\frac{Q_G}n-\frac7{50}I
-\left(\frac{\qmin(G)}n-\frac7{50}\right)I\right)
(U^{xy}u)\right]\notag\\
&=\E_{x,y}\left[\frac{A_{xy}}n
(U^{xy}u)^{\mathsf T}
\frac{Q_G-\qmin(G)I}{n}
(U^{xy}u)\right]\geq0.
\label{eq:new-matrix-rayleigh}
\end{align}
Hence
$\mathsf F-(\qmin(G)/n-7/50)\mathsf D$ is positive semidefinite.
The trace of the product of two real symmetric positive semidefinite
matrices is nonnegative~\cite[Example~2.24]{BoydVandenberghe}.
Applying this fact with the positive definite matrix $B$ and the matrix
in \eqref{eq:new-matrix-rayleigh}, and using linearity of the trace, gives
\begin{align}
0
&\leq\tr\left(B\left[
\mathsf F-\left(\frac{\qmin(G)}n-\frac7{50}\right)\mathsf D
\right]\right)\notag\\
&=\tr(B\mathsf F)
-\left(\frac{\qmin(G)}n-\frac7{50}\right)\tr(B\mathsf D).
\label{eq:new-weighted-rayleigh}
\end{align}

To divide by $\tr(B\mathsf D)$, we need its strict positivity.
Since $B$ is positive definite,
$\lambda_{\min}(B)>0$ and
$B-\lambda_{\min}(B)I$ is positive semidefinite.
Using \eqref{eq:new-D-psd} and the same trace nonnegativity gives
\begin{align*}
\tr(B\mathsf D)-\lambda_{\min}(B)\tr(\mathsf D)
&=\tr\big((B-\lambda_{\min}(B)I)\mathsf D\big)\geq0.
\end{align*}
Consequently, \eqref{eq:new-D-positive-trace} implies
\begin{equation}\label{eq:new-weighted-D-positive}
\tr(B\mathsf D)
\geq\lambda_{\min}(B)\tr(\mathsf D)>0.
\end{equation}
Rearranging \eqref{eq:new-weighted-rayleigh} and applying the hypothesis
$\tr(B\mathsf F)\leq0$ now yields
\[
\left(\frac{\qmin(G)}n-\frac7{50}\right)\tr(B\mathsf D)
\leq\tr(B\mathsf F)\leq0.
\]
Division by the positive quantity in
\eqref{eq:new-weighted-D-positive} gives
\[
\frac{\qmin(G)}n-\frac7{50}
\leq\frac{\tr(B\mathsf F)}{\tr(B\mathsf D)}\leq0,
\]
and hence $\qmin(G)\leq7n/50$.
Since $G$ was arbitrary, this proves \eqref{eq:main-bound}.
\end{proof}

In Sections~\ref{sec:definitions} and~\ref{sec:proof}, we construct a positive
definite matrix $B$ satisfying the condition in Lemma~\ref{lem:sufficient-condition}.
